\section{Localizing the fluctuation activity}\label{sec:rg-nonlinear-step}

The finite fluctuation integral in
\eqref{eq:rg-step-fluctuation-exponent} has a logarithm which is a
sum of local activities. Its construction has two localizations.
The first resolves the change of background and reconstructs
the coarse mean; the second resolves the fermion fluctuations.
After these finite operations, a connected-polymer expansion
gives the logarithm. The estimates below apply also to the first
step, where the fluctuation generators have chamber measure
$1/48$.

\subsection{Contour domains and localization}
\label{subsec:rgnl-domains}

Keep the admissible sides, mass and charge of
\cref{sec:rg-step}, and assume in addition
$\ell\ge Mr_k$. Put $e=e_k$, $h=(h_{k,1},h_{k,2})$.
The weakened photon and fermion primitives are taken on the
polydisk
\[
 \mathrm{Pd}_k=\{s=(s',s''):|s_{\mathrm{pc}}|\le M^{\alpha_0}\}.
\]
Here $s'=(s'_M,s'_{LM})$ contains the photon reconstruction
and square-root parameters, of kinds $\pi_M,\pi_{LM}$ in
\exteq{37}{MaxwellBlockManuscript}; $s''$ contains the
fermion $M$-cube parameters $\pi_k$ of
Section~\extref{4.1}{FermionBlockManuscript}.
We assume their analytic kernel and H\"older estimates,
component restriction, row vanishing and coefficient support
with common constants $(\alpha_0,\gamma_w,\kappa_w;C_w)$,
where $\kappa_w=\kappa+\kappa_0+3$.
The photon hypotheses are those of
Proposition~\extref{5.1}{MaxwellBlockManuscript},
\exteq{42}{MaxwellBlockManuscript} and the component
restriction following them. The fermion hypotheses are those of
Propositions~\extref{4.1}{FermionBlockManuscript} and
\extref{4.3}{FermionBlockManuscript}, including
\exteq{57}{FermionBlockManuscript},
\exteq{58}{FermionBlockManuscript}, and their covariance
identities. At parameter value one all families equal the
exact operators. Their gauge, lattice, transpose and reality
actions are the actions of those exact operators with the
parameters transformed accordingly. At level zero we retain
the stated diagonal families and the chamber covariance
with measure $1/48$.

The quantitative assumptions are those of \cref{thm:rgnl-current-step}.

For $\zeta\in\mathrm{Box}_k$, write
\begin{equation}\label{eq:rgnl-weakened-photon}
 Z(s')=H_k^{\mathrm{ph}}(s')C C_k^{1/2}(s')\zeta .
\end{equation}
In this formula $C$ is the real photon coordinate map of
\exteq{21}{MaxwellBlockManuscript}. The diagonal photon
family is used at $k=0$. For a bond field $Z$ set
\[
 |Z|_*=\max\{\|Z\|_\infty,\|\partial Z\|_\infty,
                              \|\delta_\alpha\partial Z\|_\infty\}.
\]
The three thresholds defining $R_k$ are
$\theta_i=e^{-3/4+i\varepsilon}$, $i=1,2,3$.
Thus $\widetilde R_k^-$ consists of the permitted real-gauge
part plus a field satisfying the three strict bounds
$(1-e^\varepsilon)\theta_i$.
Introduce the radii
\begin{equation}\label{eq:rgnl-contour-radii}
 r_1=\tfrac12e^{-3/4+9\varepsilon/2},\qquad
 r_2=e^{-3/4+7\varepsilon},\qquad
 r_4=e^{-1/4+3\varepsilon}.
\end{equation}

\begin{proposition}\label{prop:rgnl-contour-domains}
For every $s\in\mathrm{Pd}_k$ and
$\zeta\in\mathrm{Box}_k$,
\[
 |Z(s')|_*\le\tfrac12e^{-\varepsilon/2}.
\]
The radii satisfy $2\le r_2,r_4\le r_1$ and
$(r_i-1)^{-1}\le2/r_i$, $i=2,4$.
If $A\in\widetilde R_k^-$, $\mathscr C\in\mathfrak X_k$,
and $|Z|_*\le\tfrac12e^{-\varepsilon/2}$, then
\[
 A+tZ\in\widetilde R_k,\qquad
 A+\mathscr C+tZ\in\widetilde{\mathcal G}_k
                  \quad (|t|\le r_1).
\]
\end{proposition}

\begin{proof}
The square-root kernel is bounded by
$C_w\exp(-\gamma_w d(\Upsilon,\Upsilon'))$.
The summability constant $c_\gamma$ and
the coordinate bound $|\zeta(\Upsilon')|\le p_{0,k}
\le e^{-\varepsilon/4}$ give
\[
 \sup_\Upsilon|(C_k^{1/2}(s')\zeta)(\Upsilon)|
                                      \le C_wc_\gamma e^{-\varepsilon/4}.
\]
The coordinate map is real. Its bound
$\|Cv\|_\infty\le c_C\sup_\Upsilon|v(\Upsilon)|$
also gives the absolute row bound
$\sum_\Upsilon|C(b,\Upsilon)|\le c_C$:
apply the real bound to the signs of the entries of a fixed
row. Therefore
\[
 \|CC_k^{1/2}(s')\zeta\|_\infty
                                  \le c_CC_wc_\gamma e^{-\varepsilon/4}.
\]
Each of the reconstruction, derivative and H\"older derivative
kernels has the same summable exponential bound. Consequently
\[
 |Z(s')|_*\le c_CC_w^2c_\gamma^2 e^{-\varepsilon/4}
                                          \le\tfrac12 e^{-\varepsilon/2},
\]
where the last inequality is
$2c_CC_w^2c_\gamma^2e^{\varepsilon/4}\le1$.
At level zero the polynomial diagonal family has the same
row bounds: its constant is
$4\exp(2\gamma_w)M^{27\alpha_0}\le C_w$.

Since $\varepsilon\le1/200$,
\[
 \frac{r_2}{r_1}=2e^{5\varepsilon/2}\le1,\qquad
 \frac{r_4}{r_1}=2e^{1/2-3\varepsilon/2}
                                      \le2e^{5\varepsilon/2}\le1.
\]
The small-charge condition
$88M^3e^{1/4-6\varepsilon}\le1$ gives
$r_4^{-1}=e^{1/4-3\varepsilon}\le1/88$ and
$r_2^{-1}\le r_4^{-1}$.
In particular both radii are at least $2$, and their stated
inverse bounds follow.

For $|t|\le r_1$ the three norms of $tZ$ are at most
$\frac14e^{-3/4+4\varepsilon}$, strictly less than
$e^\varepsilon\theta_3$. Adding this field fills less than the
margin in each of the three defining bounds of
$\widetilde R_k^-$. The permitted gauge part is unchanged.
Hence $A+tZ\in\widetilde R_k$.
Finally $\widetilde R_k+\mathfrak X_k
\subset\widetilde{\mathcal G}_k$ gives the second inclusion.
\end{proof}
The first-step vertex estimate
\eqref{eq:rg-step-first-vertex-small-bound} now applies to
$Z(s')$ for every weakening parameter. In particular,
$96M^3e^{1/4}\le1$ gives its bound
$e^{1/4-5\varepsilon}$ in the $h/2$ norm.

Let $B=A+\mathscr C$.
The arguments of $H_k^\#$ are independent pairs
$(\Phi,\bar\Phi)$; in $H_k^\#\Phi$ the single symbol denotes
this pair. Write
$\psi^\#=(\psi,\chi;\bar\psi,\bar\chi)$ for the ordinary,
pair, barred ordinary and barred pair generators, respectively.
Then
\[
 \begin{aligned}
 H_k^\#(s'',B)\Phi
   ={}&(H_k(s'',B)\Phi,\delta_{\alpha,B}H_k(s'',B)\Phi;\\
       &
       \bar H_k(s'',B)\bar\Phi,
                       \delta_{\alpha,-B}\bar H_k(s'',B)\bar\Phi),\\
 \bar H_k(s'',B)&=\mathsf C^{-1}H_k(s'',-B)\mathsf C.
 \end{aligned}
\]
Define
$J_k^\#(s'',B,Z)=H_k^\#(s'',B+Z)-H_k^\#(s'',B)$ and
\[
 \widetilde\Gamma_k(s'',B)(W,\bar W)
                         =(\Gamma_k(s'',B)W,\bar W),\qquad
 W_k^\#(s'',B)=H_k^\#(s'',B)
                         \big(\Gamma_k(s'',B)W,\bar W\big).
\]
The notation $J_k,\bar J_k$ denotes the components without
H\"older quotient; $\mathrm{Wc},\overline{\mathrm{Wc}}$ denotes
the corresponding components of $W_k^\#$.
The auxiliary generators $W$ have weight
$\mathsf w_W=1$ for $k\ge1$ and $\sqrt{48}$ for $k=0$.

These fields are continuous in all closed polydisk and box
variables and analytic in the open parameter polydisk and the
connection variables. For the H\"older rows this follows
directly from
\[
 (\delta_{\alpha,B}H_k)((x,x'),u)
   =d_X(x,x')^{-\alpha}
       \big(U_B(x,x')H_k(x',u)-H_k(x,u)\big).
\]
The transport is entire in $B$. \Cref{prop:rgnl-contour-domains}
places every connection on the closed contours inside the
open domain of the primitives. Thus compositions at
$B+tZ(s')$ have the stated regularity also when the contour
variable is on its circle.

In the following localization, cubes and polymers are sets
of photon sites. Containment refers to all their lattice sites,
and two cubes meet when they share a lattice site.
Let $P_0(X)$ be the parameter cubes contained in $X$.
For a set $P$ of parameter cubes disjoint from $P_0(X)$, let $Y'(X,P)$ be the union of
the $M$-cubes meeting $X$ or a cube of $P_0(X)\cup P$.
For a continuous analytic family $G(s,X)$ define
\begin{align}
 \mathcal T_PG(X)
  &=\int_{[0,1]^P}\dd s_P\,
       \partial_PG(s,X)
       \big|_{s=1\text{ on }P_0(X),\,
                    s=0\text{ outside }P_0(X)\cup P},
                          \label{eq:rgnl-interpolation}\\
 \operatorname{Loc}G(Y)
  &=\sum_{(X,P):Y'(X,P)=Y}\mathcal T_PG(X).
                          \label{eq:rgnl-localization}
\end{align}
Repeated one-variable integration gives the useful identity
\begin{equation}\label{eq:rgnl-finite-difference}
 \mathcal T_PG(X)=
    \sum_{Q\subseteq P}(-1)^{|P|-|Q|}
                                 G(\mathbf1_{P_0(X)\cup Q},X).
\end{equation}
In particular, coefficient locality can be checked at
indicator families; no statement about the convergence of
the later logarithm is used here.

\begin{lemma}\label{lem:rgnl-resummation}
Suppose
$\|G(s,X)\|\le B_0\exp(-\kappa'd_M(X))$ on the parameter
polydisk, where $\kappa_0\le\kappa'\le\kappa$.
Assume $\mathcal T_PG(X)=0$ whenever $Y'(X,P)$ is
disconnected. Then
\begin{gather}
 \sum_Y\operatorname{Loc}G(Y)=\sum_XG(1,X),
                         \label{eq:rgnl-localization-sum}\\
 \|\operatorname{Loc}G(Y)\|
   \le4c_XK_0B_0
             \exp\big(-(\kappa'-\kappa_0-1)d_M(Y)\big).
                         \label{eq:rgnl-resummation-bound}
\end{gather}
\end{lemma}

\begin{proof}
Applying $f(1)=f(0)+\int_0^1f'(t)\dd t$ in each parameter
outside $P_0(X)$ gives $G(1,X)=\sum_P\mathcal T_PG(X)$.
The sums are finite. After removing the vanishing disconnected
terms, each term belongs to exactly one polymer $Y'(X,P)$.
This proves \eqref{eq:rgnl-localization-sum}.

Put $R=M^{\alpha_0}$. Cauchy's estimate on a polydisk of
radius less than $R-1$ about each integration point, followed
by the limit of these radii to $R-1$, gives
\[
 \|\mathcal T_PG(X)\|
       \le B_0\exp(-\kappa'd_M(X))(R-1)^{-|P|}.
\]
If $Y'(X,P)=Y$, then $X\subseteq Y$ and
$|Y\setminus X|_M\le N_c|P|$, since every added $M$-cube
meets a parameter cube and each parameter cube meets at
most $N_c=(2L+1)^3$ such cubes.
To compare their tree lengths, order the cubes of $Y\setminus X$
by graph distance from $X$. Each added cube is adjacent
to a preceding one. Extend a net by a segment from
a point of that preceding closed cube to their common
face, of length at most $M$. Starting from an
arbitrarily near-minimal net for $X$ gives
$d_M(Y)\le d_M(X)+|Y\setminus X|_M$.
Conversely every net for $Y$ is a net for $X$.
Thus $0\le d_M(Y)-d_M(X)\le N_c|P|$.
Writing $a_X=\log(6c_X)$, the condition
$R-1\ge\exp(N_c\kappa+a_X)$ implies
\[
 \|\mathcal T_PG(X)\|
 \le B_0\exp(-(\kappa'-\kappa_0)d_M(Y))
             \exp(-\kappa_0d_M(X))\exp(-a_X|P|).
\]
At most $3|Y|_M$ parameter cubes are contained in $Y$.
Hence
\[
 \sum_P\exp(-a_X|P|)
       \le(1+(6c_X)^{-1})^{3|Y|_M}
       \le \exp((d_M(Y)+1)/2).
\]
The polymer counting and size estimates give
\[
 \sum_{X\subseteq Y}\exp(-\kappa_0d_M(X))
   \le K_0|Y|_M
   \le2c_XK_0\exp(d_M(Y)/2).
\]
The product of the last two bounds is at most
$4c_XK_0\exp(d_M(Y))$. Summing the estimate for
$\mathcal T_PG(X)$ proves
\eqref{eq:rgnl-resummation-bound}.
\end{proof}

\subsection{The two finite localizations}
\label{subsec:rgnl-finite-localizations}

For a function analytic near the closed disk $|t|\le r$, $r>1$,
write
\[
 \mathcal C_rf=\frac1{2\pi i}\int_{|t|=r}
                                      \frac{f(t)}{t(t-1)}\,\dd t .
\]
This is $f(1)-f(0)$ and satisfies
$\|\mathcal C_rf\|\le(r-1)^{-1}\sup_{|t|=r}\|f(t)\|$.
Let $(m,\mathscr E)\in\operatorname{Ball}_k^x(4)$ and set
$m'=m+m(\mathscr E)$, $\mathscr E^r=R\mathscr E$.
For $k\ge1$ define
\begin{align}
 G_1(s,X)&=\mathcal C_{r_1}
                \langle\bar\Psi,M_k(X,B+tZ(s'))\Psi\rangle_1,
                                      \label{eq:rgnl-G1}\\
 G_2(s,X)&=\mathcal C_{r_2}\,
       \mathscr E^r(X;A+tZ(s'),\mathscr C,
                    \psi^\#+tJ_k^\#(s'',B,Z(s'))\Psi),
                                      \label{eq:rgnl-G2}\\
 G_3(s,\square)&=\mathcal C_{r_2}(-m')
       \langle\bar\psi+t\bar J_k(s'',B,Z(s'))\bar\Psi,
                         \psi+tJ_k(s'',B,Z(s'))\Psi\rangle_\square .
                                      \label{eq:rgnl-G3}
\end{align}
The symbols $M_k(X,\cdot)$ denote the local Dirac primitives.
A function defined for one cube is zero on the other polymers.
Put
\[
 E_i^*(Y)=\operatorname{Loc}G_i(Y),\qquad
 E_i(Y)=\operatorname{sub}_{Y,B}E_i^*(Y),\qquad
 E_{\le3}=\sum_{i=1}^3E_i.
\]
The substitution is \eqref{eq:rg-step-local-substitution}.
At $k=0$ there is no $\Psi$ and no substitution:
$G_1(s,\square)=\mathcal V_0(Z(s'),\square)$,
$G_2=G_3=0$, and $E_{\le3}=\operatorname{Loc}G_1$.
All parameter kinds occur in this first localization.

For the second localization, which uses only $\pi_k$, let
$v=W_k^\#(s'',B)$ and define
\begin{align}
 G_{\le3}^{\,W}(s'',X)&=E_{\le3}(X;A,\mathscr C,\zeta,
                                                      \psi^\#+v),
                                      \label{eq:rgnl-GW}\\
 G_4(s'',X)&=\mathcal C_{r_4}
                  \mathscr E^r(X;A,\mathscr C,\psi^\#+tv),\\
 G_5(s'',\square)&=\mathcal C_{r_4}(-m')
                 \langle\bar\psi+t\overline{\mathrm{Wc}},
                                   \psi+t\mathrm{Wc}\rangle_\square,\\
 F(Y)&=\operatorname{Loc}(G_{\le3}^{\,W}+G_4+G_5)(Y).
                                      \label{eq:rgnl-finite-activity}
\end{align}
At $k=0$ the term $G_4$ is zero and $m'=m$.

\begin{proposition}\label{prop:rgnl-finite-locality}
All finite expressions just defined exist on
$\widetilde R_k^-\times\mathfrak X_k\times\mathrm{Box}_k$.
The element $F(Y)$ is even, has ordinary legs over $Y$ and
pair legs admissible for $Y$, and has fluctuation generators
only at unit indices in $Y$. Its coefficients read $A,\mathscr C$
only on bonds meeting $Y$ and read $\zeta$ only through the
coordinates $\zeta(\Upsilon)$ with $\supp\Upsilon\subset Y$.
They are continuous in these arguments and analytic in the
connection variables. The same assertions hold for $E_{\le3}(Y)$,
which has no fluctuation generators.

The disconnected terms in both localizations vanish, and
\begin{equation}\label{eq:rgnl-two-localization-sums}
 \sum_YE_i^*(Y)=\sum_XG_i(1,X),\qquad
 \sum_YF(Y)=\sum_X(G_{\le3}^{\,W}+G_4+G_5)(1,X).
\end{equation}
The finite expressions on $Y$ depend on the datum only through
$m'$ and the $\mathscr E^r(X)$ with $X\subseteq Y$.
\end{proposition}

\begin{proof}
We give the component argument which is responsible for
locality. Let $\mathrm{Ps}=P_0(X)\cup P$, let
$Y'=Y'(X,P)$ and let $Y'_1$ be its component containing $X$.
Set the weakening parameters to zero outside $\mathrm{Ps}$.
For each of the three parameter kinds, any link-component
$\mathrm{Lnk}$ of $\mathrm{Ps}$ is contained in the polymer
formed by the $M$-cubes meeting it. This polymer lies in one
component of $Y'$. If a site of $Y'_1$ lies in $\mathrm{Lnk}$,
or in the enlarged parameter cube of a photon kind, that
polymer has a cube in common with $Y'_1$, and hence lies
entirely in $Y'_1$.

The row-vanishing property of the weakened photon families
and their component restriction give, on any fine bond $b$
meeting $Y'_1$,
\begin{equation}\label{eq:rgnl-restricted-photon}
 Z(s')(b)=
  \sum_{\substack{b',\Upsilon,\Upsilon'\subset Y'_1}}
   H_k^{\mathrm{ph}}(s')(b,b')C(b',\Upsilon)
          C_k^{1/2}(s')(\Upsilon,\Upsilon')\zeta(\Upsilon').
\end{equation}
Indeed, the reconstruction row is zero unless $b$ lies in
one link-component; its columns lie in that component.
The coordinate map has columns supported in the $L$-blocks
of the ends of $b'$, and each block lies in the $M$-cube of
the corresponding end. The square-root row then has columns
in its own link-component. Each component involved is
attached to $Y'_1$ by the preceding argument.
These rows are unchanged if the parameters outside $Y'_1$
are changed while their common zero set outside $\mathrm{Ps}$
is kept.

For a chamber site $x=(c,g)$ at distance at most $1$ from
$Y'_1$, put $K_c=\{c':|c'-c|_\infty\le2L\}$, as in
\cite[Proposition~4.1]{FermionBlockManuscript}.
The fermion row is zero unless every parameter cube
meeting $K_c$ is present. These cubes form one
link-component. When the row is nonzero, that component
contains $c$ and a site of $Y'_1$, so it lies in $Y'_1$.
The row of $H_k(s'',B)$, and its barred counterpart,
then has columns in that component and reads $B$ only on
bonds meeting it. At a column of this row, the corresponding
row of $\Gamma_k(s'',B)$ has columns in the same component.
These conclusions also hold at level zero: the diagonal
factor is
\[
 \prod_{\mathrm{pc}\in\pi_0:\mathrm{pc}\cap K_c\ne\varnothing}
                                                       s''_{\mathrm{pc}},
\]
and the kernel of that multiplication operator is
$48\mathbf1_{x=x'}$ times this factor. Its composition with
$\Gamma_0$ uses the measure $1/48$.

An endpoint of a pair admissible for $X$ is at distance
at most $\eta_k\le1$ from $X$. Its transport reads only
bonds contained in $X$, or the one type-$3$ bond meeting $X$.
Thus the preceding row argument applies to every row
used in $J_k^\#$ and $W_k^\#$.
The H\"older row has only the two columns $x,x'$ and
the entire transport coefficient displayed above.
Finite compositions preserve this dependence: the columns
of $O_2O_1$ are the union of the columns of $O_1$ at the
finitely many columns of $O_2$, and its read bonds are the
union of their read bonds. Entries of the transport and
reconstruction rows are entire in the connections, and finite
compositions of these algebraic rows retain this property.
Compositions involving weakened primitives are holomorphic
on their stated primitive domains. These are the finite-row
properties used here; the weakened operators themselves
are not asserted to have finite range.

For $G_1$ at $k\ge1$, the local primitive $M_k(X,\cdot)$
has both indices in $X$ and reads only bonds meeting $X$.
For $G_2,G_3$, the incoming activity or mass form has ordinary
legs in $X$ and admissible pair legs; the substitution by
$J_k^\#\Psi$ introduces $\Psi$ only in $Y'_1$.
At $k=0$ the vertex has precisely these ordinary and
type-$3$ pair legs and reads $Z$ only on bonds meeting
its cube. Equation~\eqref{eq:rgnl-restricted-photon}
therefore proves the same component dependence for every
first-localization function.

If $Y'$ is disconnected, choose an $M$-cube outside $Y'_1$.
It cannot meet $X$ or a parameter cube contained in $X$,
so it meets a cube of $P$ outside $Y'_1$. The integrand is
independent of that parameter. Pairing the subsets in
\eqref{eq:rgnl-finite-difference} according to whether they
contain this cube proves that the term is zero.
For a term assigned to a connected $Y$, all parameter cubes
in $P_0(X)\cup P$ and $X$ lie in $Y$.
Equation~\eqref{eq:rgnl-finite-difference} now gives the
claimed locality of $E_i^*$.

The reconstruction substitution adds only ordinary legs
over $Y$ and boundary pairs $(x,x^{(3)})$ admissible for
$Y$. Its finite rows at a unit site $y$ read bonds contained
in $\Delta_y$ or type-$3$ bonds from this block. These bonds
meet $Y$ since $\Delta_y\subset Y$. The barred rows have
the same support. Thus it preserves the claimed coefficient
locality and removes $\Psi$.
For the second localization, the row argument for
$H_k^\#\widetilde\Gamma_kW$ gives fluctuation generators
only in $Y'_1$. It applies to each of the three functions
in \eqref{eq:rgnl-finite-activity}. The same disconnected
parameter cancellation proves the second assertion.
\Cref{lem:rgnl-resummation} gives both sum identities.

Regularity follows from the established contour domains,
the primitive regularity, and finite sums and substitutions.
Contour integrals preserve analyticity: their Riemann sums
converge uniformly on compact sets. The same argument
applies to finite parameter integrals. The incoming
activities and mass forms are even, and each substitution
sends generators to odd linear combinations, so every
expression remains even.
Finally $G_1$ is independent of the datum; $G_2,G_4$
use $\mathscr E^r(X)$ and $G_3,G_5$ use $m'$.
Tracing the two localizations shows that no incoming
polymer outside $Y$ enters the result.
\end{proof}

The proof identifies the primitive entries determining a
finite expression on $Y$. With parameters zero outside $Y$,
the photon field reads only
$H_k^{\mathrm{ph}}(s')(b,b')$ for $b$ meeting $Y$,
$b'\subset Y$, and
$C_k^{1/2}(s')(\Upsilon,\Upsilon')$ with both supports
in $Y$. Its remaining factor is the fixed coordinate map.
The fermion fields use $H_k,\bar H_k$ at chamber sites
over $Y$, with unit columns in $Y$, at $B+Z(s')$ and
$B$, and $\Gamma_k(s'',B)$ with both unit indices in $Y$.
At zero, the kernel sums have measure $1/48$ and the
diagonal rows have the factor $48$ already recorded.
The Dirac primitive has both indices in $X\subseteq Y$
and connection $B+tZ(s')$ on $|t|=r_1$.
H\"older transports read only bonds meeting $Y$.
The reconstruction substitution contains only the fixed
finite $T_k^{(i)},\bar T_k^{(i)}$ rows at unit sites of its
own subpolymer.

Consequently two systems with the stated hypotheses on this torus give the
same finite expressions on $Y$ if these entries agree
for every parameter family zero outside $Y$, at the
connections formed with their respective photon fields.
The restricted photon sum first gives equality of the
fields on bonds meeting $Y$; the transports and incoming
activities there then agree by their locality.
The listed fermion and Dirac entries give equal first
functions, their localizations and reconstruction
substitutions. They give equal second functions and
localizations in turn, and then equal activities $K(Y)$ and
$H(Y)$ in \eqref{eq:rgnl-Mayer-activity} and
\eqref{eq:rgnl-integrated-activity}, and each term of the
connected series \eqref{eq:rgnl-log-activity}.
Agreement of the photon
entries and agreement of the fermion entries as functions
on the full operator domain is sufficient: if the two
composed connections differ outside $Y$, local row
dependence removes that difference first.
The same comparison passes to the convergent logarithm
once its convergence has been established.

\subsection{Bounds before fluctuation integration}
\label{subsec:rgnl-finite-bounds}

Set
\[
 \begin{gathered}
 \kappa_w=\kappa+\kappa_0+3,\qquad
 a_1=\kappa-\kappa_0-1,\qquad a_2=\kappa-2\kappa_0-2,\\
 C_W=28C_w^2c_\gamma,
 \qquad H_{\mathrm{cl}}=36c_X^2K_0^2e^{1/4-5\varepsilon}.
 \end{gathered}
\]
The notation $H_{\mathrm{cl}}$ is a scalar cluster bound. For a
scalar-component kernel $O(a,b)$ with column measure $\nu(b)$, write
$|O|_{1,\infty}=\sup_a\sum_b\nu(b)|O(a,b)|$, including spin indices.
This is the row norm of \cite[Theorem~B]{NormalizedContraction}.
The finite primitive bounds imply
\begin{align}
 \|\langle\bar\Psi,M_k(X,B)\Psi\rangle_1\|_{w_3}
 &\le4c_Xc_\gamma C_wM^3w_3^2
                   \exp(-(\kappa_w-1)d_M(X)),\quad k\ge1,
                                      \label{eq:rgnl-M-bound}\\
 |H_k^\#(s'',B)|_{1,\infty}&\le4C_w,\\
 |J_k^\#(s'',B,Z)|_{1,\infty}
                  &\le16C_we^{3/4-9\varepsilon/2},\\
 \mathsf w_W|H_k^\#(s'',B)\widetilde\Gamma_k(s'',B)|_{1,\infty}
                  &\le C_W.           \label{eq:rgnl-W-bound}
\end{align}
Here $B=A+\mathscr C$ and $|Z|_*\le\frac12e^{-\varepsilon/2}$.
To see the first bound, count $4M^3|X|_M$ rows and sum
the exponential kernel in its second index, giving
$c_\gamma C_w \exp(-\kappa_wd_M(X))$ per row.
Then $|X|_M\le c_X\exp(d_M(X))$ gives
\eqref{eq:rgnl-M-bound}.
For $H_k^\#$ the primitive bound controls the ordinary
and short-pair rows by $C_w$. A long pair, $d_X\ge1$,
has its H\"older quotient bounded by four times the
ordinary sup norm. At level zero every pair is long.
The monomial matrices $\mathsf C^{\pm1}$ preserve row norms,
so the barred rows have the same bounds.
Apply $\mathcal C_{r_1}$ to
$H_k^\#(s'',B+tZ)$ to obtain the bound for $J_k^\#$.
Finally $|\Gamma_k|_{1,\infty}\le C_wc_\gamma$ by kernel
summation, using $1/48$ at level zero. Submultiplicativity
and $\mathsf w_W\le7$ give \eqref{eq:rgnl-W-bound}.

\begin{proposition}\label{prop:rgnl-finite-bounds}
For the data of \cref{prop:rgnl-finite-locality},
\begin{align}
 \|E_{\le3}(Y)\|_{h/2}
       &\le8c_XK_0e^{1/4-5\varepsilon}\exp(-a_1d_M(Y)),
                                      \label{eq:rgnl-E3-bound}\\
 \|F(Y)\|_{(h/4;\mathsf w_W)}
       &\le H_{\mathrm{cl}}\exp(-a_2d_M(Y)).
                                      \label{eq:rgnl-F-bound}
\end{align}
The element $F(Y)$ is the sum of a term independent of
$(m,\mathscr E)$ and a term complex linear in these data.
\end{proposition}

\begin{proof}
Extraction gives
\[
 \|\mathscr E^r(X)\|_h\le2\|\mathscr E\|_k^x\exp(-\kappa d_M(X))
                   \le12e^{1/4-7\varepsilon}\exp(-\kappa d_M(X)),
 \qquad |m'|\le8e^{3/4-8\varepsilon}.
\]
The mass bound uses
$c_{\mathrm{fl}}e^\varepsilon\le1$.
With chamber measure $\eta_k^3/48$, direct counting gives
\[
 \|\langle\bar\psi,\psi\rangle_\square\|_{(w_1,w_2)}
                                                    =4M^3w_1^2.
\]
The number of monomials is
$4\cdot48\eta_k^{-3}M^3$ and their common coefficient
is $\eta_k^3/48$.

In the first localization give $\Psi$ the weight
$w_3=C_Th_{k,1}$ and the old generators the weights $h/2$.
The contour bound and \eqref{eq:rgnl-M-bound} give
\[
 \|G_1(s,X)\|\le
 16c_Xc_\gamma C_wC_T^2M^3e^{1/4-9\varepsilon/2}
                 \exp(-(\kappa_w-1)d_M(X))
 \le e^{1/4-5\varepsilon}\exp(-\kappa d_M(X)).
\]
The last step uses
\[
 16c_Xc_\gamma C_wC_T^2M^3e^{\varepsilon/2}\le1.
\]
At $k=0$ the weighted vertex instead gives
$12M^3e^{1/2-\varepsilon/2}\le e^{1/4-5\varepsilon}$.
Moreover,
\[
 r_2|J_k^\#|_{1,\infty}w_3
     \le16C_wC_Te^{3\varepsilon/2}h_{k,2}
     \le\tfrac12h_{k,2}.
\]
The substitution bound therefore takes the integrand norm
of $G_2$ to the $h$ norm of $\mathscr E^r$, and that of
$G_3$ to the $h$ norm of the mass form. It gives
\[
 \|G_2(s,X)\|\le24e^{1-14\varepsilon}\exp(-\kappa d_M(X)),
 \qquad \|G_3(s,\square)\|\le64M^3e^{1-15\varepsilon}.
\]
The condition $88M^3e^{1/4-6\varepsilon}\le1$ implies
\[
 24e^{3/4-9\varepsilon}
                   +64M^3e^{3/4-10\varepsilon}\le1.
\]
Thus $\sum_{i\le3}G_i$ is bounded by
$2e^{1/4-5\varepsilon}\exp(-\kappa d_M(X))$.
Apply \cref{lem:rgnl-resummation}, using the disconnected
cancellation in \cref{prop:rgnl-finite-locality}.
The reconstruction substitution costs no norm, since
$C_T(h_{k,1}+h_{k,2})/2\le w_3$.
This proves \eqref{eq:rgnl-E3-bound}.

For the second localization, the condition
$112C_w^2c_\gamma e^{2\varepsilon}\le3$ gives
\[
 r_4C_W\le\tfrac34h_{k,2},\qquad C_W\le\tfrac14h_{k,2}.
\]
Consequently addition of $W_k^\#$ maps the $h/4$ norm
to the $h/2$ norm at $t=1$ and to the $h$ norm on
the contour $|t|=r_4$. The three input bounds are
\begin{align*}
 \|G_{\le3}^{\,W}(s'',X)\|
       &\le8c_XK_0e^{1/4-5\varepsilon}\exp(-a_1d_M(X)),\\
 \|G_4(s'',X)\|&\le24e^{1/2-10\varepsilon}\exp(-\kappa d_M(X)),\\
 \|G_5(s'',\square)\|&\le64M^3e^{1/2-11\varepsilon}.
\end{align*}
Their sum is at most
$9c_XK_0e^{1/4-5\varepsilon}\exp(-a_1d_M(X))$,
again by $88M^3e^{1/4-6\varepsilon}\le1$.
Since $\kappa_0\le a_1\le\kappa$, resummation gives
\eqref{eq:rgnl-F-bound}.
At level zero $G_4=0$; all the other inequalities include
the factor $\mathsf w_W$ already, so they apply unchanged.
Linearity follows from the definitions: $G_1$ is fixed,
$G_2,G_4$ are linear in $R\mathscr E$, $G_3,G_5$ in $m'$,
and localization, the reconstruction substitution and the
contour functionals are linear.
\end{proof}

